\subsection{Evaluation Metrics}
\label{sec:evaluation_metrics}

We evaluate the generated PAMAP2 windows from two complementary perspectives:
similarity to the real sensor-window distribution and downstream human-activity
recognition (HAR) utility. This follows the evaluation logic used by
SDForger~\cite{rousseau2025sdforger}, which separates similarity metrics from
task-utility metrics, and the TSGBench taxonomy~\cite{ang2023tsgbench}, which
groups time-series similarity metrics into feature-based and distance-based
measures. Lower values are better for all similarity metrics, while higher
values are better for HAR utility metrics.

For each activity class $a$, let
$\mathcal{R}_a=\{X_i\}_{i=1}^{n_a}$ denote sampled real windows and
$\mathcal{G}_a=\{\tilde{X}_i\}_{i=1}^{m_a}$ denote sampled generated windows,
where each window $X_i,\tilde{X}_i\in\mathbb{R}^{L\times C}$ has length $L$ and
$C$ sensor channels. The implementation samples at most a fixed number of
windows per activity, computes an activity-wise score, and then reports the
macro-average over activities.

\paragraph{Feature-based similarity.}
The first group measures whether generated windows preserve real-data
statistics and temporal structure. Marginal Distribution Difference (MDD)
compares empirical marginal distributions over time steps and channels:
\begin{equation}
    \mathrm{MDD}(a)=
    \frac{1}{LCB}
    \sum_{t=1}^{L}\sum_{c=1}^{C}\sum_{b=1}^{B}
    \left|h_{t,c}^{R}(b)-h_{t,c}^{G}(b)\right| ,
\end{equation}
where $h_{t,c}^{R}$ and $h_{t,c}^{G}$ are real and generated histogram density
estimates and $B$ is the number of bins.
Autocorrelation Difference (ACD) compares the lag-wise autocorrelation of the
real and generated windows:
\begin{equation}
    \mathrm{ACD}(a)=
    \frac{1}{C}
    \sum_{c=1}^{C}
    \left\|\boldsymbol{\rho}_{c}^{R}
    -\boldsymbol{\rho}_{c}^{G}\right\|_2 ,
\end{equation}
where $\boldsymbol{\rho}_{c}$ is the autocorrelation vector for channel $c$.
Skewness Difference (SD) compares the distribution asymmetry of real
and generated windows:
\begin{equation}
    \mathrm{SD}(a)=
    \frac{1}{C}
    \sum_{c=1}^{C}
    \left|
    \gamma(r_{a,c})-\gamma(g_{a,c})
    \right| ,
\end{equation}
where $\gamma(\cdot)$ denotes skewness. Finally, Kurtosis Difference (KD)
compares the excess kurtosis of real and generated channel distributions:
\begin{equation}
    \mathrm{KD}(a)=
    \frac{1}{C}
    \sum_{c=1}^{C}
    \left|
    \kappa(r_{a,c})-\kappa(g_{a,c})
    \right| ,
\end{equation}
where $r_{a,c}$ and $g_{a,c}$ are the flattened real and generated values for
channel $c$, and $\kappa(\cdot)$ denotes excess kurtosis.

\paragraph{Distance-based similarity.}
The second group compares paired real and generated windows directly. Euclidean
Distance (ED) is averaged over paired samples and channels:
\begin{equation}
    \mathrm{ED}(a)=
    \frac{1}{N_a C}
    \sum_{i=1}^{N_a}\sum_{c=1}^{C}
    \left\|X_{i,:,c}-\tilde{X}_{i,:,c}\right\|_2 .
\end{equation}
Dynamic Time Warping (DTW) measures temporally flexible alignment between paired
real and generated windows~\cite{sakoe1978dtw}:
\begin{equation}
    \mathrm{DTW}(a)=
    \frac{1}{N_a}
    \sum_{i=1}^{N_a}
    D_{\mathrm{DTW}}\!\left(X_i,\tilde{X}_i\right),
\end{equation}
where
\begin{equation}
    D_{\mathrm{DTW}}(u,v)=
    \min_{\pi}
    \sum_{(p,q)\in\pi}\left\|u_p-v_q\right\|_2
\end{equation}
over monotonic warping paths $\pi$. We also compute the SDForger shapelet
reconstruction metric SHR/SHAP-RE when enabled. Shapelets are learned from real
windows and then used to reconstruct generated windows:
\begin{equation}
    \mathrm{SHR}(a)
    =
    \frac{1}{N_a}
    \sum_{i=1}^{N_a}
    \frac{1}{2}
    \left\|
    \tilde{X}_i-\alpha_i\mathcal{T}_{\delta_i}(S)
    \right\|_2^2 ,
\end{equation}
where $S$ is the learned shapelet basis, $\alpha_i$ are sparse coefficients, and
$\mathcal{T}_{\delta_i}$ shifts the basis by offsets $\delta_i$. When
available, we additionally report ED/real and DTW/real, obtained by dividing the
real--synthetic score by the corresponding real-vs-real baseline. These ratios
are interpreted within a fixed dataset, channel subset, window length, and
normalization protocol.

\paragraph{Downstream HAR utility.}
Similarity alone does not guarantee that synthetic data is useful for a
downstream task. We therefore evaluate HAR utility with a
RandomForestClassifier~\cite{breiman2001randomforest}. Each window is flattened
into a feature vector and the classifier is evaluated on the held-out real test
set under three training settings:
\begin{equation}
    \mathcal{D}_{\mathrm{train}}=
    \begin{cases}
    \mathcal{D}_{\mathrm{real}}, & \text{real-only},\\
    \mathcal{D}_{\mathrm{syn}}, & \text{synthetic-only},\\
    \mathcal{D}_{\mathrm{real}}\cup\mathcal{D}_{\mathrm{syn}},
    & \text{real-plus-synthetic}.
    \end{cases}
\end{equation}
The real-only setting is the supervised baseline, synthetic-only corresponds to
train-on-synthetic-test-on-real utility, and real-plus-synthetic measures whether
generated windows improve or degrade augmentation performance. For each setting
we report accuracy,
\begin{equation}
    \mathrm{Accuracy}=
    \frac{1}{N}\sum_{i=1}^{N}\mathbf{1}[\hat{y}_i=y_i],
\end{equation}
and macro-F1,
\begin{equation}
    \mathrm{MacroF1}=
    \frac{1}{|\mathcal{Y}|}
    \sum_{y\in\mathcal{Y}}
    \frac{2\,\mathrm{Precision}_y\,\mathrm{Recall}_y}
    {\mathrm{Precision}_y+\mathrm{Recall}_y}.
\end{equation}
Confusion matrices are additionally used to inspect class-wise failure modes.

\paragraph{Implementation note.}
The metrics above describe the exact evaluation used in our current repository:
MDD, ACD, SD, KD, ED, DTW, optional SHR, and Random-Forest HAR utility with
real-only, synthetic-only, and real-plus-synthetic settings. The abbreviation SD
is therefore written explicitly as Skewness Difference, matching the
TSGBench/SDForger metric description~\cite{ang2023tsgbench,rousseau2025sdforger}.

% BibTeX entries used by this section:
%
% @article{ang2023tsgbench,
%   title        = {{TSGBench}: Time Series Generation Benchmark},
%   author       = {Ang, Yihao and Huang, Qiang and Bao, Yifan and Tung, Anthony K. H. and Huang, Zhiyong},
%   journal      = {Proceedings of the VLDB Endowment},
%   volume       = {17},
%   number       = {3},
%   pages        = {305--318},
%   year         = {2023}
% }
%
% @inproceedings{rousseau2025sdforger,
%   title     = {Forging Time Series with Language: A Large Language Model Approach to Synthetic Data Generation},
%   author    = {Rousseau, C{\'e}cile and Boschi, Tobia and Cornacchia, Giandomenico and Salwala, Dhaval and Pascale, Alessandra and Moreno, Juan Bernabe},
%   booktitle = {Advances in Neural Information Processing Systems},
%   year      = {2025}
% }
%
% @article{sakoe1978dtw,
%   title   = {Dynamic Programming Algorithm Optimization for Spoken Word Recognition},
%   author  = {Sakoe, Hiroaki and Chiba, Seibi},
%   journal = {IEEE Transactions on Acoustics, Speech, and Signal Processing},
%   volume  = {26},
%   number  = {1},
%   pages   = {43--49},
%   year    = {1978}
% }
%
% @article{breiman2001randomforest,
%   title   = {Random Forests},
%   author  = {Breiman, Leo},
%   journal = {Machine Learning},
%   volume  = {45},
%   number  = {1},
%   pages   = {5--32},
%   year    = {2001}
% }
